% Dokumen "serba ada": memakai paket yang paling sering dipakai paper, supaya
% semuanya ikut ke dalam cache yang dibundel. Tidak untuk dibaca siapa pun.
\documentclass[11pt,a4paper]{article}
\usepackage[T1]{fontenc}
\usepackage[utf8]{inputenc}
\usepackage[bahasa,english]{babel}
\usepackage{lmodern}
\usepackage{amsmath,amssymb,amsthm,mathtools}
\usepackage{graphicx,float,wrapfig}
\usepackage{booktabs,tabularx,longtable,multirow,array,colortbl}
\usepackage{siunitx}
\usepackage{caption,subcaption}
\usepackage{enumitem}
\usepackage{geometry}
\usepackage{setspace}
\usepackage{fancyhdr}
\usepackage{titlesec}
\usepackage{xcolor}
\usepackage{hyperref}
\usepackage{url}
\usepackage{listings}
\usepackage{algorithm}
\usepackage{algpseudocode}
\usepackage{tikz}
\usetikzlibrary{arrows.meta,positioning,shapes.geometric,calc,fit,backgrounds,patterns}
\usepackage{pgfplots}
\pgfplotsset{compat=1.18}
\usepackage{natbib}
\usepackage{csquotes}

\title{Serba Ada}
\author{Pengumpul Paket}

\begin{document}
\maketitle
\tableofcontents

\section{Matematika}
\begin{equation}
  E = mc^2, \qquad \int_0^\infty e^{-x^2}\,dx = \frac{\sqrt{\pi}}{2}.
\end{equation}
\begin{align}
  a &= b + c \\
  d &\le \sum_{i=1}^{n} x_i
\end{align}

\section{Tabel}
\begin{table}[H]
  \centering
  \caption{Contoh}
  \begin{tabular}{lrr}
    \toprule
    Metode & \si{\percent} & Waktu (\si{\second}) \\
    \midrule
    Dasar & 91 & 12 \\
    \bottomrule
  \end{tabular}
\end{table}

\section{Gambar}
\begin{tikzpicture}
  \node[draw,rounded corners] (a) {Mulai};
  \node[draw,right=of a] (b) {Selesai};
  \draw[-Stealth] (a) -- (b);
\end{tikzpicture}

\begin{tikzpicture}
  \begin{axis}[width=6cm,height=4cm]
    \addplot coordinates {(0,0) (1,1) (2,4)};
  \end{axis}
\end{tikzpicture}

\section{Algoritma}
\begin{algorithm}
\caption{Contoh}
\begin{algorithmic}
\State $x \gets 0$
\end{algorithmic}
\end{algorithm}

\section{Kode}
\begin{lstlisting}[language=Python]
def halo():
    return "dunia"
\end{lstlisting}

\end{document}
